\section{Grafi}

\subsection{Rappresentazione dei grafi}

\subsubsection{Definizione}
In questa parte del corso si assume che i concetti matematici di base sui grafi siano gia' noti. Si richiamano pero' le due rappresentazioni fondamentali usate dagli algoritmi:
\begin{itemize}
	\item \textbf{matrice di adiacenza};
	\item \textbf{liste di adiacenza}.
\end{itemize}

Sia
\[
G=(V,E)
\]
un grafo, orientato o non orientato. Indichiamo con
\[
n=|V|,\qquad m=|E|
\]
il numero di vertici e di archi.

\paragraph{Matrice di adiacenza.}
Dato un grafo di $n$ nodi, la matrice di adiacenza e' una matrice $n\times n$ in cui righe e colonne sono etichettate con i nomi dei nodi. La cella di coordinate $(i,j)$ indica se esiste l'arco dal nodo $i$ al nodo $j$:
\[
A[i,j]=
\begin{cases}
1 & \text{se } (i,j)\in E,\\
0 & \text{altrimenti.}
\end{cases}
\]
Nel caso di grafo non orientato, la matrice e' simmetrica.

\paragraph{Liste di adiacenza.}
La rappresentazione con liste di adiacenza mantiene una lista principale dei nodi e, per ciascun nodo $v$, una lista
\[
Adj[v]
\]
contenente tutti i vertici $u$ tali che
\[
(v,u)\in E.
\]
Se il numero di nodi e' fisso, la lista principale puo' essere implementata come un array di riferimenti a liste.

\subsubsection{Intuizione}
La matrice di adiacenza risponde molto velocemente alla domanda:
\[
\text{``l'arco }(u,v)\text{ esiste?''}
\]
perche' basta leggere una cella. Tuttavia usa sempre spazio quadratico, anche se il grafo ha pochi archi.

Le liste di adiacenza, invece, memorizzano soltanto gli archi presenti. Sono quindi piu' adatte per grafi sparsi, cioe' grafi in cui $m$ e' molto piu' piccolo di $n^2$.

\subsubsection{Motivazione}
La scelta della rappresentazione incide direttamente sulla complessita' degli algoritmi. BFS, DFS, ordinamento topologico e componenti fortemente connesse vengono analizzati nelle slide con liste di adiacenza, perche' in tale rappresentazione l'esplorazione di tutti gli archi costa
\[
\Theta(|V|+|E|).
\]

\subsubsection{Proprieta' teoriche}
\[
\begin{array}{c|c|c}
\text{Operazione} & \text{Matrice di adiacenza} & \text{Liste di adiacenza} \\
\toprule
\text{Memoria} & \Theta(n^2) & \Theta(n+m) \\
\text{Verificare se }(u,v)\in E & \Theta(1) & \Theta(d(u)) \\
\text{Accedere agli adiacenti di }u & \Theta(n) & \Theta(d(u)) \\
\text{Identificare tutti gli archi} & \Theta(n^2) & \Theta(n+m) \\
\bottomrule
\end{array}
\]

Nel caso di grafi non orientati, ogni arco $\{u,v\}$ compare due volte nelle liste di adiacenza: una volta in $Adj[u]$ e una volta in $Adj[v]$.

\subsubsection{Algoritmo}
Per costruire una matrice di adiacenza da una lista di archi:
\begin{algorithm}[H]
	\caption{\textsc{CostruisciMatriceAdiacenza}$(V,E)$}
	\KwIn{Insieme di vertici $V=\{1,\ldots,n\}$, insieme di archi $E$}
	\KwOut{Matrice di adiacenza $A$}
	\For{$i\gets 1$ \KwTo $n$}{
		\For{$j\gets 1$ \KwTo $n$}{
			$A[i,j]\gets 0$\;
		}
	}
	\ForEach{$(u,v)\in E$}{
		$A[u,v]\gets 1$\;
	}
	\Return $A$\;
\end{algorithm}

Per un grafo non orientato, quando si incontra l'arco $\{u,v\}$ si impostano entrambe le celle:
\[
A[u,v]\gets 1,\qquad A[v,u]\gets 1.
\]

\subsubsection{Esempio guidato}
Consideriamo il grafo non orientato
\[
V=\{1,2,4,5\},\qquad
E=\{\{1,2\},\{1,4\},\{2,4\},\{4,5\}\}.
\]
Con ordine dei vertici $1,2,4,5$, la matrice di adiacenza e':
\[
\begin{array}{c|cccc}
 & 1 & 2 & 4 & 5 \\
\toprule
1 & 0 & 1 & 1 & 0 \\
2 & 1 & 0 & 1 & 0 \\
4 & 1 & 1 & 0 & 1 \\
5 & 0 & 0 & 1 & 0 \\
\bottomrule
\end{array}
\]
La stessa informazione con liste di adiacenza e':
\[
\begin{array}{c|l}
1 & 2,4 \\
2 & 1,4 \\
4 & 1,2,5 \\
5 & 4
\end{array}
\]

\subsubsection{Grafico o diagramma}
\begin{figure}[H]
	\centering
	\begin{tikzpicture}[
		v/.style={circle, draw=primary, thick, minimum size=7mm, font=\small},
		e/.style={primary, thick}
		]
		\node[v] (n1) at (0,1.2) {1};
		\node[v] (n2) at (1.6,1.2) {2};
		\node[v] (n4) at (0.8,0) {4};
		\node[v] (n5) at (2.4,0) {5};
		\draw[e] (n1)--(n2) (n1)--(n4) (n2)--(n4) (n4)--(n5);
	\end{tikzpicture}
	\hspace{1.5cm}
	\raisebox{1.2cm}{\begin{minipage}{0.38\textwidth}
		\footnotesize
		\textbf{Liste di adiacenza:}
		\begin{itemize}[leftmargin=*, nosep, topsep=2pt]
			\item $\mathrm{Adj}[1]=\langle 2,4\rangle$
			\item $\mathrm{Adj}[2]=\langle 1,4\rangle$
			\item $\mathrm{Adj}[4]=\langle 1,2,5\rangle$
			\item $\mathrm{Adj}[5]=\langle 4\rangle$
		\end{itemize}
	\end{minipage}}
	\caption{Grafo non orientato d'esempio e liste di adiacenza.}
\end{figure}

\subsubsection{Dimostrazione}
Le slide non presentano una dimostrazione formale delle rappresentazioni; forniscono definizioni e costi. I costi seguono direttamente dalla struttura dati:
\begin{itemize}
	\item una matrice $n\times n$ contiene sempre $n^2$ celle;
	\item una lista di adiacenza contiene un riferimento per ogni nodo e un elemento di lista per ogni arco orientato;
	\item in un grafo non orientato ogni arco compare in due liste.
\end{itemize}

\subsubsection{Analisi della complessita'}
La matrice richiede memoria $\Theta(n^2)$, indipendentemente dal numero di archi.

Le liste richiedono memoria $\Theta(n+m)$, perche' servono $n$ riferimenti principali e spazio proporzionale agli archi memorizzati.

Se si accetta di usare piu' memoria del necessario, e' possibile implementare le liste di adiacenza in modo da ridurre quasi a tempo costante la verifica di esistenza di un nodo o di un arco. Le slide pongono questa osservazione come domanda; il contenuto implementativo dettagliato non viene sviluppato nel materiale ufficiale.

\subsubsection{Errori tipici d'esame}
\begin{hardnote}
Confondere il costo di verificare un arco con il costo di scandire tutti gli adiacenti. Nella matrice il primo e' $\Theta(1)$, il secondo e' $\Theta(n)$.
\end{hardnote}

\begin{hardnote}
Dimenticare che, nei grafi non orientati rappresentati con liste di adiacenza, ogni arco compare due volte.
\end{hardnote}

\subsubsection{Possibili domande d'esame}
\begin{itemize}
	\item Definire matrice di adiacenza e liste di adiacenza.
	\item Confrontare i costi delle due rappresentazioni.
	\item Spiegare perche' la matrice di un grafo non orientato e' simmetrica.
	\item Dire quale rappresentazione conviene per grafi sparsi.
\end{itemize}

\subsubsection{Collegamenti teorici}
BFS, DFS, ordinamento topologico e componenti fortemente connesse sono analizzati sulle liste di adiacenza. Per questo la complessita' tipica sara'
\[
\Theta(|V|+|E|).
\]

% ===============================================================
\subsection{Visita in ampiezza}

\subsubsection{Definizione}
La \textbf{visita in ampiezza}, o \textbf{BFS} (\emph{breadth-first search}), e' un algoritmo di ricerca su grafi.

Dato un grafo
\[
G=(V,E)
\]
e un nodo sorgente
\[
s\in V,
\]
BFS ispeziona gli archi di $G$ per scoprire tutti i nodi raggiungibili da $s$.

La visita procede per livelli: per
\[
k=1,2,\ldots,n-1,
\]
scopre i nodi a distanza $k$ a partire dai nodi a distanza $k-1$.

L'algoritmo:
\begin{itemize}
	\item calcola la distanza minima da $s$ a ciascun nodo raggiungibile;
	\item genera un albero, detto \textbf{albero BF} o \textbf{breadth-first tree}, con radice $s$ e contenente tutti i nodi raggiungibili da $s$.
\end{itemize}

\subsubsection{Intuizione}
BFS usa una coda FIFO. La coda contiene i nodi scoperti ma non ancora completamente esaminati. Questo garantisce che i nodi scoperti prima vengano espansi prima.

Di conseguenza, prima vengono visitati tutti i nodi a distanza $0$, poi quelli a distanza $1$, poi quelli a distanza $2$, e cosi' via. Questa e' la ragione per cui BFS calcola cammini minimi nei grafi non pesati.

\subsubsection{Motivazione}
BFS e' uno schema di base per molti algoritmi importanti sui grafi. Le slide citano esplicitamente:
\begin{itemize}
	\item algoritmo di Prim per l'albero di connessione minimo;
	\item algoritmo di Dijkstra per i cammini minimi da sorgente unica.
\end{itemize}

In questa lezione BFS viene usata soprattutto per mostrare come si esplora un grafo per livelli e come si dimostra formalmente la correttezza delle distanze calcolate.

\subsubsection{Proprieta' teoriche}
Ogni nodo ha tre attributi:
\[
u.color,\qquad u.d,\qquad u.\pi.
\]

\begin{itemize}
	\item $u.color$ indica lo stato del nodo.
	\item $u.d$ e' la distanza stimata da $s$.
	\item $u.\pi$ e' il predecessore di $u$ nel cammino da $s$ a $u$.
\end{itemize}

I colori hanno il seguente significato:
\[
\begin{array}{c|l}
\text{Colore} & \text{Significato} \\
\toprule
\textsc{White} & \text{nodo non ancora scoperto} \\
\textsc{Gray} & \text{nodo scoperto ma non ancora completato} \\
\textsc{Black} & \text{nodo completato: tutti i suoi vicini sono stati scoperti} \\
\bottomrule
\end{array}
\]

L'invariante di ciclo indicato nelle slide per il ciclo \texttt{while} e':
\[
\text{quando viene eseguito il test del ciclo \texttt{while}, la coda }Q
\text{ contiene solo nodi grigi.}
\]

\subsubsection{Algoritmo}
\begin{algorithm}[H]
	\caption{\textsc{BFS}$(G,s)$}
	\KwIn{$G=(V,E)$, sorgente $s\in V$}
	\ForEach{$u\in V\setminus\{s\}$}{
		$u.color\gets\textsc{White}$\;
		$u.d\gets\infty$\tcp*{$d$ e' la distanza da $s$}
		$u.\pi\gets\textsc{Null}$\tcp*{$\pi$ e' il predecessore}
	}
	$s.color\gets\textsc{Gray}$\;
	$s.d\gets 0$\;
	$s.\pi\gets\textsc{Null}$\;
	$Q\gets\emptyset$\tcp*{coda FIFO dei nodi grigi}
	$Q.\textsc{Enqueue}(s)$\;
	\While{$Q\neq\emptyset$}{
		$u\gets Q.\textsc{Dequeue}()$\;
		\ForEach{$v\in G.Adj[u]$}{
			\If{$v.color=\textsc{White}$}{
				$v.color\gets\textsc{Gray}$\;
				$v.d\gets u.d+1$\;
				$v.\pi\gets u$\;
				$Q.\textsc{Enqueue}(v)$\;
			}
		}
		$u.color\gets\textsc{Black}$\;
	}
\end{algorithm}

\subsubsection{Esempio guidato}
Usiamo l'esempio coerente con la visita mostrata nelle slide, con sorgente $2$.
Le liste di adiacenza rilevanti sono:
\[
\begin{array}{c|l}
1 & 2,4 \\
2 & 1,4 \\
3 & 6 \\
4 & 2,1,5 \\
5 & 4 \\
6 & 3
\end{array}
\]

La componente raggiungibile da $2$ contiene i nodi $2,1,4,5$. I nodi $3$ e $6$ non sono raggiungibili da $2$.

\begin{center}
\begin{tabular}{c|c|c|c|c}
\toprule
Passo & Operazione & Coda $Q$ & Distanze finite & Predecessori nuovi \\
\midrule
$0$ & inizializza $s=2$ & $\langle 2\rangle$ & $2.d=0$ & -- \\
$1$ & esamina $2$, scopre $1$ & $\langle 1\rangle$ & $1.d=1$ & $1.\pi=2$ \\
$2$ & esamina $2$, scopre $4$ & $\langle 1,4\rangle$ & $4.d=1$ & $4.\pi=2$ \\
$3$ & completa $2$ & $\langle 1,4\rangle$ & -- & -- \\
$4$ & esamina $1$ & $\langle 4\rangle$ & -- & -- \\
$5$ & esamina $4$, scopre $5$ & $\langle 5\rangle$ & $5.d=2$ & $5.\pi=4$ \\
$6$ & completa $5$ & $\langle\rangle$ & -- & -- \\
\bottomrule
\end{tabular}
\end{center}

Al termine:
\[
2.d=0,\qquad 1.d=1,\qquad 4.d=1,\qquad 5.d=2,
\]
mentre
\[
3.d=6.d=\infty.
\]

L'albero BF e':
\[
2\to 1,\qquad 2\to 4,\qquad 4\to 5.
\]

\subsubsection{Grafico o diagramma}
\begin{figure}[H]
	\centering
	\begin{tikzpicture}[
		v/.style={circle, draw=primary, thick, minimum size=7mm, font=\small},
		w/.style={v, fill=white},
		g/.style={v, fill=gray!35},
		b/.style={v, fill=primary!85, text=white},
		e/.style={-{Stealth[length=2mm]}, primary, thick}
		]
		\node[w] (n1) at (0,1.2) {1};
		\node[g] (n2) at (1.6,1.2) {2};
		\node[g] (n4) at (0.8,0) {4};
		\node[w] (n5) at (2.4,0) {5};
		\node[w] (n3) at (4.2,1.2) {3};
		\node[w] (n6) at (4.2,0) {6};
		\draw[e] (n1)--(n2) (n1)--(n4) (n2)--(n4) (n4)--(n5) (n3)--(n6);
		\node[font=\scriptsize, below=2mm of n2] {sorgente $s=2$};
		\node[font=\scriptsize, anchor=west] at (5.0,1.5) {
			\begin{tabular}{@{}l@{}}
				\cellcolor{white}\rule{0pt}{2.5mm}\;White\; non scoperto\\
				\cellcolor{gray!35}\;Gray\; in coda\\
				\cellcolor{primary!85}\textcolor{white}{\;Black\; completato}
			\end{tabular}
		};
	\end{tikzpicture}
	\caption{\textsc{BFS}$(G,2)$ dopo la scansione del nodo $2$: $1,4$ grigi in coda, $2$ nero.}
\end{figure}

\begin{figure}[H]
	\centering
	\begin{tikzpicture}[
		v/.style={circle, draw=primary, thick, minimum size=7mm, font=\small},
		w/.style={v, fill=white},
		e/.style={primary, thick},
		tree/.style={-{Stealth[length=2mm]}, primary, thick},
		lbl/.style={font=\scriptsize, primary}
		]
		\node[w] (n1) at (0,1.2) {1};
		\node[w] (n2) at (1.6,1.2) {2};
		\node[w] (n4) at (0.8,0) {4};
		\node[w] (n5) at (2.4,0) {5};
		\node[w] (n3) at (4.2,1.2) {3};
		\node[w] (n6) at (4.2,0) {6};
		\draw[e, gray!50] (n1)--(n2) (n1)--(n4) (n2)--(n4) (n4)--(n5) (n3)--(n6);
		\node[lbl, above=1mm of n2] {$d{=}0$};
		\node[lbl, above=1mm of n1] {$d{=}1$};
		\node[lbl, below=1mm of n4] {$d{=}1$};
		\node[lbl, below=1mm of n5] {$d{=}2$};
		\node[lbl, above=1mm of n3] {$d{=}\infty$};
		\node[lbl, below=1mm of n6] {$d{=}\infty$};
	\end{tikzpicture}
	\hspace{1.2cm}
	\begin{tikzpicture}[
		v/.style={circle, draw=primary, thick, minimum size=7mm, font=\small},
		tree/.style={-{Stealth[length=2mm]}, primary, thick}
		]
		\node[v] (r) at (0,1.4) {2};
		\node[v] (l) at (-1.0,0) {1};
		\node[v] (m) at (0.2,0) {4};
		\node[v] (d) at (1.4,-1.0) {5};
		\draw[tree] (r)--(l) (r)--(m) (m)--(d);
		\node[font=\scriptsize, below=2mm of d] {albero BF};
	\end{tikzpicture}
	\caption{Grafo con distanze finali da $s{=}2$ (sinistra) e albero BF con archi $\pi$ (destra).}
\end{figure}

\subsubsection{Dimostrazione}
La correttezza di BFS rispetto ai cammini minimi nei grafi non pesati viene dimostrata nelle slide tramite tre lemmi, un corollario e un teorema.

\begin{definizione}[Grafo pesato]
Un grafo $G=(V,E)$ si dice pesato quando esiste una funzione peso
\[
\omega:E\to\mathbb{R}
\]
che associa a ogni arco un valore detto peso o valore dell'arco. Un grafo non pesato puo' essere visto come un grafo pesato in cui
\[
\omega(e)=1\qquad \forall e\in E.
\]
\end{definizione}

\begin{definizione}[Distanza di cammino minimo]
Si definisce distanza di cammino minimo da $s$ a $v$ il valore
\[
\delta(s,v)=
\min_{P\,:\,s\overset{P}{\leadsto} v}
\sum_{e\in P}\omega(e).
\]
Se non esiste un cammino da $s$ a $v$, allora
\[
\delta(s,v)=\infty.
\]
Un cammino da $s$ a $v$ di lunghezza $\delta(s,v)$ e' detto cammino minimo.
\end{definizione}

\begin{lemma}
Se $G=(V,E)$ e' un grafo non pesato e $s\in V$, allora per qualsiasi arco $(u,v)\in E$ vale
\[
\delta(s,v)\le \delta(s,u)+1.
\]
\end{lemma}

\begin{proof}
Se $u$ e' raggiungibile da $s$, sia
\[
s\overset{P'}{\leadsto} u
\]
un cammino minimo da $s$ a $u$. Allora concatenando $P'$ con l'arco $(u,v)$ si ottiene un cammino da $s$ a $v$ di lunghezza
\[
\delta(s,u)+1.
\]
Il cammino minimo da $s$ a $v$ non puo' essere piu' lungo di un cammino specifico da $s$ a $v$, quindi
\[
\delta(s,v)\le \delta(s,u)+1.
\]
Se $u$ non e' raggiungibile da $s$, allora $\delta(s,u)=\infty$ e la disuguaglianza e' ancora vera.
\end{proof}

\begin{lemma}
Sia $G=(V,E)$ un grafo non pesato e sia $s\in V$. Se si esegue \textsc{BFS}$(G,s)$, al termine vale
\[
v.d\ge \delta(s,v)\qquad\forall v\in V.
\]
\end{lemma}

\begin{proof}
La dimostrazione e' per induzione sul numero di operazioni \textsc{Enqueue}.

\textbf{Caso base.}
All'inizio solo $s$ e' inserito in $Q$:
\[
s.d=0=\delta(s,s).
\]
Per ogni altro nodo $v$,
\[
v.d=\infty\ge \delta(s,v).
\]

\textbf{Passo induttivo.}
Supponiamo che la proprieta' valga prima di una nuova operazione \textsc{Enqueue}. Un nodo bianco $v$ viene scoperto durante la visita di un nodo $u$. Per ipotesi induttiva:
\[
u.d\ge \delta(s,u).
\]
L'algoritmo assegna:
\[
v.d\gets u.d+1.
\]
Quindi:
\[
v.d=u.d+1\ge \delta(s,u)+1\ge \delta(s,v),
\]
dove l'ultima disuguaglianza segue dal lemma precedente. Dopo questa assegnazione, $v.d$ non viene piu' modificato, perche' $v$ non sara' piu' bianco.
\end{proof}

\begin{lemma}
Si assuma che durante BFS la coda contenga
\[
Q=\langle v_1,v_2,\ldots,v_r\rangle.
\]
Allora:
\[
v_r.d\le v_1.d+1
\]
e
\[
v_i.d\le v_{i+1}.d\qquad\text{per }i=1,2,\ldots,r-1.
\]
\end{lemma}

\begin{proof}
La dimostrazione e' per induzione sulle operazioni possibili sulla coda $Q$.

\textbf{Base.}
Quando
\[
Q=\langle s\rangle,
\]
la tesi e' immediata.

\textbf{Passo induttivo: estrazione.}
Se si estrae $v_1$, la coda diventa
\[
\langle v_2,\ldots,v_r\rangle.
\]
Per ipotesi induttiva valevano
\[
v_r.d\le v_1.d+1
\]
e
\[
v_1.d\le v_2.d.
\]
Da queste due relazioni segue
\[
v_r.d\le v_2.d+1.
\]
L'ordine non decrescente delle distanze nella coda rimane quindi valido.

\textbf{Passo induttivo: inserimento.}
Se durante l'esame di $u=v_1$ si inserisce un nuovo nodo $v$ in coda, allora
\[
v.d=u.d+1=v_1.d+1.
\]
Dopo l'estrazione di $u$, il nuovo primo elemento della coda ha distanza almeno $v_1.d$. Pertanto il nuovo ultimo elemento soddisfa ancora la proprieta' di stare a distanza al piu' uno dal primo elemento della coda.
\end{proof}

\begin{corollario}
Se $v_i$ e $v_j$ sono inseriti in coda durante l'esecuzione di BFS, e $v_i$ viene inserito prima di $v_j$, allora
\[
v_i.d\le v_j.d.
\]
\end{corollario}

\begin{teorema}[Correttezza di BFS]
\textsc{BFS} scopre tutti i nodi $v\in V$ raggiungibili da $s$ e, alla fine,
\[
v.d=\delta(s,v)\qquad\forall v\in V.
\]
Inoltre, per qualsiasi nodo $v$ raggiungibile da $s$, un cammino minimo da $s$ a $v$ si ottiene prendendo un cammino minimo da $s$ a $v.\pi$ e aggiungendo l'arco $(v.\pi,v)$.
\end{teorema}

\begin{proof}
La dimostrazione procede per assurdo, come nelle slide.

Supponiamo che esista almeno un vertice raggiungibile per cui la distanza calcolata non e' minima. Tra tutti questi vertici, scegliamo $v$ con distanza minima $\delta(s,v)$.

Dal lemma precedente sappiamo che
\[
v.d\ge \delta(s,v).
\]
Poiche' per ipotesi $v.d\neq \delta(s,v)$, deve valere
\[
v.d>\delta(s,v).
\]
Il nodo $v$ e' diverso da $s$, perche'
\[
s.d=0=\delta(s,s).
\]

Sia $u$ il vertice che precede immediatamente $v$ in un cammino minimo
\[
s\overset{P}{\leadsto}v.
\]
Allora:
\[
\delta(s,v)=\delta(s,u)+1.
\]
Poiche' $u$ precede $v$ nel cammino minimo, vale
\[
\delta(s,u)<\delta(s,v).
\]
Per la scelta di $v$ come vertice errato di minima distanza, $u$ non puo' essere errato. Quindi:
\[
u.d=\delta(s,u).
\]
Ne segue:
\[
v.d>\delta(s,v)=\delta(s,u)+1=u.d+1.
\]

Consideriamo il momento in cui BFS esamina $v$ come adiacente di $u$.

\textbf{Caso 1: $v$ e' bianco.}
L'algoritmo assegna
\[
v.d\gets u.d+1,
\]
in contraddizione con
\[
v.d>u.d+1.
\]

\textbf{Caso 2: $v$ e' nero.}
Allora $v$ e' gia' stato rimosso dalla coda prima di $u$. Per il corollario sull'ordine di inserimento e sulle distanze:
\[
v.d\le u.d.
\]
Quindi:
\[
v.d\le u.d<u.d+1<v.d,
\]
contraddizione.

\textbf{Caso 3: $v$ e' grigio.}
Allora $v$ e' stato colorato grigio durante la visita di un nodo $w$, rimosso dalla coda prima di $u$, e
\[
v.d=w.d+1.
\]
Per il corollario:
\[
w.d\le u.d.
\]
Quindi:
\[
v.d=w.d+1\le u.d+1,
\]
in contraddizione con
\[
v.d>u.d+1.
\]

Tutti i casi portano a contraddizione. Dunque per ogni nodo raggiungibile $v$ vale
\[
v.d=\delta(s,v).
\]

Quando un nodo $v$ viene scoperto, l'algoritmo assegna $v.\pi=u$ e $v.d=u.d+1$. Per quanto dimostrato, $u.d=\delta(s,u)$ e $v.d=\delta(s,v)$, quindi il cammino minimo verso $v$ si ottiene ricorsivamente dal cammino minimo verso $v.\pi$ seguito dall'arco $(v.\pi,v)$.
\end{proof}

\subsubsection{Analisi della complessita'}
Il primo ciclo \texttt{for} e le inizializzazioni della sorgente preparano i nodi:
\[
O(|V|).
\]
Le slide precisano che si scrive $O$ e non necessariamente $\Theta$, perche' l'inizializzazione puo' essere effettuata in modo pigro in $O(1)$.

Il ciclo \texttt{while} continua finche' ci sono nodi grigi. Ogni nodo diventa grigio una sola volta, quindi viene inserito in coda una sola volta e rimosso una sola volta. Le operazioni \textsc{Enqueue} e \textsc{Dequeue} costano $O(1)$.

Per ciascun nodo $u$, il ciclo sugli adiacenti esamina tutta la lista $Adj[u]$. Sommando su tutti i nodi:
\[
\sum_{u\in V}|Adj[u]|=\Theta(|E|)
\]
nei grafi orientati; nei grafi non orientati ogni arco e' contato due volte, ma il costo resta $\Theta(|E|)$.

La complessita' finale e':
\[
O(|V|+|E|).
\]
Lo spazio aggiuntivo e' $O(|V|)$ per colori, distanze, predecessori e coda.

\subsubsection{Errori tipici d'esame}
\begin{hardnote}
Dire che BFS calcola cammini minimi in qualunque grafo pesato. Nelle slide la correttezza dei cammini minimi e' dimostrata per grafi non pesati, equivalenti a grafi con peso unitario su ogni arco.
\end{hardnote}

\begin{hardnote}
Confondere nodo grigio e nodo nero. Un nodo grigio e' stato scoperto ma non completato; un nodo nero ha gia' avuto tutti i vicini esaminati.
\end{hardnote}

\begin{hardnote}
Dimenticare che la coda e' FIFO. Senza FIFO, l'esplorazione per livelli e la dimostrazione sulle distanze non valgono.
\end{hardnote}

\subsubsection{Possibili domande d'esame}
\begin{itemize}
	\item Scrivere lo pseudocodice di BFS.
	\item Spiegare il ruolo dei colori.
	\item Dimostrare l'invariante sulla coda dei nodi grigi.
	\item Dimostrare che BFS calcola le distanze minime nei grafi non pesati.
	\item Simulare BFS mostrando coda, colori, distanze e predecessori.
	\item Analizzare la complessita' $O(|V|+|E|)$.
\end{itemize}

\subsubsection{Collegamenti teorici}
BFS e' la base per ragionare su:
\begin{itemize}
	\item cammini minimi in grafi non pesati;
	\item costruzione di alberi di visita;
	\item algoritmi che usano code o code di priorita' per espandere nodi in un ordine controllato.
\end{itemize}

% ===============================================================
\subsection{Visita in profondita'}

\subsubsection{Definizione}
La \textbf{visita in profondita'}, o \textbf{DFS} (\emph{depth-first search}), visita un grafo andando il piu' possibile in profondita'.

Dato un grafo
\[
G=(V,E),
\]
DFS:
\begin{itemize}
	\item ispeziona gli archi a partire dall'ultimo nodo scoperto che ha ancora archi non ispezionati;
	\item quando tutti gli archi di un nodo $v$ sono stati ispezionati, ritorna al nodo da cui $v$ e' stato scoperto;
	\item continua finche' esistono nodi da ispezionare.
\end{itemize}

La visita puo' generare piu' alberi: il risultato e' una \textbf{foresta DF}, cioe' una \emph{depth-first forest}.

\subsubsection{Intuizione}
DFS si comporta come una pila ricorsiva. Appena scopre un nuovo nodo, sospende temporaneamente l'esame del nodo corrente e scende sul nuovo nodo. Solo quando non puo' piu' scendere, torna indietro.

Questa differenza e' centrale:
\[
\begin{array}{c|c}
\text{BFS} & \text{DFS} \\
\toprule
\text{esplora per livelli} & \text{esplora andando in profondita'} \\
\text{usa una coda FIFO} & \text{usa ricorsione, cioe' una pila implicita} \\
\text{calcola distanze minime non pesate} & \text{calcola tempi di scoperta e completamento} \\
\bottomrule
\end{array}
\]

\subsubsection{Motivazione}
DFS e' usata nelle slide come base per:
\begin{itemize}
	\item classificazione degli archi;
	\item proprieta' a parentesi dei tempi;
	\item ordinamento topologico;
	\item componenti fortemente connesse.
\end{itemize}

\subsubsection{Proprieta' teoriche}
Per ogni nodo $v$, DFS mantiene:
\[
v.color,\qquad v.\pi,\qquad v.d,\qquad v.f.
\]

\begin{itemize}
	\item $v.d$ e' l'istante in cui $v$ viene scoperto, cioe' diventa grigio.
	\item $v.f$ e' l'istante in cui $v$ viene completato, cioe' diventa nero.
	\item $v.\pi$ e' il predecessore di $v$ nella foresta DF.
\end{itemize}

L'algoritmo usa una variabile globale
\[
time
\]
che viene incrementata ogni volta che un nodo cambia colore per scoperta o completamento.

\paragraph{Classificazione degli archi.}
Le slide riportano la seguente classificazione degli archi ispezionati da DFS:
\[
\begin{array}{c|l}
\text{Tipo} & \text{Condizione} \\
\toprule
T\ \text{(tree)} & \text{l'arco }(u,v)\text{ e' ispezionato e }v\text{ e' bianco} \\
B\ \text{(backward)} & \text{l'arco }(u,v)\text{ e' ispezionato e }v\text{ e' grigio} \\
F\ \text{(forward)} & \text{l'arco }(u,v)\text{ e' ispezionato, }v\text{ e' nero e }u.d<v.d \\
C\ \text{(cross)} & \text{l'arco }(u,v)\text{ e' ispezionato, }v\text{ e' nero e }u.d>v.f \\
\bottomrule
\end{array}
\]

Per un arco backward, il nodo $v$ e' un antenato di $u$. Per un arco forward, $v$ e' un discendente di $u$. Per un arco cross, $u$ e $v$ non sono in relazione di discendenza, oppure appartengono a due alberi distinti della foresta DF.

\subsubsection{Algoritmo}
\begin{algorithm}[H]
	\caption{\textsc{DFS}$(G)$}
	\KwIn{$G=(V,E)$}
	\ForEach{$u\in V$}{
		$u.color\gets\textsc{White}$\;
		$u.\pi\gets\textsc{Null}$\;
	}
	$time\gets 0$\tcp*{variabile globale}
	\ForEach{$u\in V$}{
		\If{$u.color=\textsc{White}$}{
			\textsc{DFS-Visit}$(G,u)$\;
		}
	}
\end{algorithm}

\begin{algorithm}[H]
	\caption{\textsc{DFS-Visit}$(G,u)$}
	\KwIn{$G=(V,E)$, nodo $u\in V$}
	$time\gets time+1$\;
	$u.d\gets time$\;
	$u.color\gets\textsc{Gray}$\;
	\ForEach{$v\in G.Adj[u]$}{
		\If{$v.color=\textsc{White}$}{
			$v.\pi\gets u$\;
			\textsc{DFS-Visit}$(G,v)$\;
		}
	}
	$time\gets time+1$\;
	$u.f\gets time$\;
	$u.color\gets\textsc{Black}$\;
\end{algorithm}

\subsubsection{Esempio guidato}
L'esecuzione mostrata nelle slide produce la foresta:
\[
2\to 1\to 4\to 5,
\qquad
3\to 6.
\]
I tempi sono:
\[
\begin{array}{c|cc}
\text{Nodo} & d & f \\
\toprule
2 & 1 & 8 \\
1 & 2 & 7 \\
4 & 3 & 6 \\
5 & 4 & 5 \\
3 & 9 & 12 \\
6 & 10 & 11 \\
\bottomrule
\end{array}
\]

Evoluzione della pila ricorsiva:
\[
\begin{array}{c|c|c}
\text{Evento} & time & \text{Stack chiamate DFS-Visit} \\
\toprule
\text{scopri }2 & 1 & \langle 2\rangle \\
\text{scopri }1 & 2 & \langle 1,2\rangle \\
\text{scopri }4 & 3 & \langle 4,1,2\rangle \\
\text{scopri }5 & 4 & \langle 5,4,1,2\rangle \\
\text{completa }5 & 5 & \langle 5,4,1,2\rangle \\
\text{completa }4 & 6 & \langle 4,1,2\rangle \\
\text{completa }1 & 7 & \langle 1,2\rangle \\
\text{completa }2 & 8 & \langle 2\rangle \\
\text{scopri }3 & 9 & \langle 3\rangle \\
\text{scopri }6 & 10 & \langle 6,3\rangle \\
\text{completa }6 & 11 & \langle 6,3\rangle \\
\text{completa }3 & 12 & \langle 3\rangle \\
\bottomrule
\end{array}
\]

\subsubsection{Grafico o diagramma}
\begin{figure}[H]
	\centering
	\begin{tikzpicture}[
		v/.style={circle, draw=primary, thick, minimum size=7mm, font=\small},
		e/.style={-{Stealth[length=2mm]}, primary, thick},
		lbl/.style={font=\scriptsize, primary}
		]
		\node[v] (n2) at (0,2.4) {2};
		\node[v] (n1) at (0,1.2) {1};
		\node[v] (n4) at (0,0) {4};
		\node[v] (n5) at (0,-1.2) {5};
		\node[v] (n3) at (2.8,1.2) {3};
		\node[v] (n6) at (2.8,0) {6};
		\draw[e] (n2)--(n1) (n1)--(n4) (n4)--(n5) (n3)--(n6);
		\node[lbl, right=1mm of n2] {$d{=}1,\ f{=}8$};
		\node[lbl, right=1mm of n1] {$d{=}2,\ f{=}7$};
		\node[lbl, right=1mm of n4] {$d{=}3,\ f{=}6$};
		\node[lbl, right=1mm of n5] {$d{=}4,\ f{=}5$};
		\node[lbl, right=1mm of n3] {$d{=}9,\ f{=}12$};
		\node[lbl, right=1mm of n6] {$d{=}10,\ f{=}11$};
	\end{tikzpicture}
	\hspace{1.2cm}
	\begin{minipage}{0.42\textwidth}
		\centering
		\footnotesize
		\textbf{Espressione a parentesi:}\\[0.3em]
		$(2\ (1\ (4\ (5\ 5)\ 4)\ 1)\ 2)$\\[0.4em]
		$(3\ (6\ 6)\ 3)$
	\end{minipage}
	\caption{Foresta DF dell'esempio e tempi $d/f$ (schema delle slide).}
\end{figure}

\subsubsection{Dimostrazione}
Le slide non richiedono le dimostrazioni dei teoremi 20.7, 20.9 e 20.10. Si esplicita pero' la proprieta' a parentesi indicata nel materiale.

\paragraph{Struttura a parentesi.}
La scoperta di un nodo $u$ si rappresenta con una parentesi aperta $(u$, mentre il completamento di $u$ si rappresenta con una parentesi chiusa $u)$.

Ogni chiamata ricorsiva \textsc{DFS-Visit}$(G,u)$:
\begin{itemize}
	\item apre la parentesi di $u$ quando assegna $u.d$;
	\item completa interamente tutte le chiamate ricorsive dei discendenti;
	\item chiude la parentesi di $u$ quando assegna $u.f$.
\end{itemize}
Quindi non e' possibile chiudere un nodo prima di aver chiuso tutti i suoi discendenti. L'espressione prodotta e' ben formata.

Nell'esempio delle slide:
\[
(2\ (1\ (4\ (5\ 5)\ 4)\ 1)\ 2)(3\ (6\ 6)\ 3)
\]
mostra esattamente che l'intervallo temporale di un discendente e' contenuto nell'intervallo temporale del suo antenato.

\subsubsection{Analisi della complessita'}
L'algoritmo \textsc{DFS} richiede tempo $\Theta(|V|)$ escluso il tempo delle chiamate a \textsc{DFS-Visit}.

\textsc{DFS-Visit} viene eseguito esattamente una volta per ciascun nodo, quindi ci sono $|V|$ chiamate.

In ogni chiamata \textsc{DFS-Visit}$(G,u)$, il ciclo \texttt{foreach} viene eseguito
\[
|Adj[u]|
\]
volte. Sommando su tutti i nodi:
\[
\sum_{u\in V}|Adj[u]|=\Theta(|E|).
\]

Il costo totale e':
\[
\Theta(|V|+|E|).
\]
Lo spazio aggiuntivo e' $O(|V|)$ per colori, predecessori, tempi e stack ricorsivo.

\subsubsection{Errori tipici d'esame}
\begin{hardnote}
Confondere $v.d$ con una distanza. In DFS, $v.d$ e' un tempo di scoperta, non una distanza dalla sorgente.
\end{hardnote}

\begin{hardnote}
Pensare che DFS produca sempre un solo albero. Se il grafo non e' tutto raggiungibile dal primo nodo esaminato, DFS produce una foresta.
\end{hardnote}

\begin{hardnote}
Dimenticare che l'ordine di ispezione dei vertici puo' cambiare la foresta DF e i tempi, anche se le applicazioni dell'algoritmo restano corrette.
\end{hardnote}

\subsubsection{Possibili domande d'esame}
\begin{itemize}
	\item Scrivere \textsc{DFS} e \textsc{DFS-Visit}.
	\item Simulare DFS mostrando stack ricorsivo, colori, tempi $d$ e $f$.
	\item Spiegare la differenza tra BFS e DFS.
	\item Definire archi tree, backward, forward e cross.
	\item Spiegare la struttura a parentesi.
	\item Analizzare la complessita' $\Theta(|V|+|E|)$.
\end{itemize}

\subsubsection{Collegamenti teorici}
DFS e' la base per:
\begin{itemize}
	\item ordinamento topologico dei DAG;
	\item componenti fortemente connesse;
	\item ragionamenti sui tempi di scoperta e completamento.
\end{itemize}

% ===============================================================
\subsection{Ordinamento topologico}

\subsubsection{Definizione}
Un \textbf{DAG} e' un grafo diretto aciclico, cioe' un \emph{directed acyclic graph}.

Un \textbf{ordinamento topologico} di un DAG
\[
G=(V,E)
\]
e' un ordinamento lineare dei nodi tale che, per ogni arco
\[
(u,v)\in E,
\]
il nodo $u$ compare prima del nodo $v$ nell'ordinamento.

\subsubsection{Intuizione}
Un ordinamento topologico dispone i nodi su una linea in modo che tutti gli archi vadano da sinistra verso destra.

Se un arco $(u,v)$ rappresenta il vincolo ``$u$ deve avvenire prima di $v$'', allora un ordinamento topologico e' una lista che rispetta tutte le precedenze.

\subsubsection{Motivazione}
Le slide indicano che gli ordinamenti topologici sono usati in molti contesti per rappresentare precedenze fra eventi. L'algoritmo proposto usa i tempi di completamento calcolati da DFS.

\subsubsection{Proprieta' teoriche}
La proprieta' fondamentale riportata nelle slide e':
\[
\text{in un DAG, per ogni arco }(u,v)\in E,\quad u.f>v.f.
\]
Quindi, se i nodi sono inseriti in testa a una lista al momento del completamento, il nodo con tempo di completamento maggiore finisce prima nella lista.

Poiche' un DAG non contiene cicli, durante DFS non possono comparire archi all'indietro che chiuderebbero un ciclo.

\subsubsection{Algoritmo}
\begin{algorithm}[H]
	\caption{\textsc{Topological-Sort}$(G)$}
	\KwIn{DAG $G=(V,E)$}
	\KwOut{Un ordinamento topologico dei nodi}
	$lista\gets\langle\rangle$\;
	Esegui \textsc{DFS}$(G)$ e, ogni volta che un valore $u.f$ viene calcolato, inserisci $u$ in testa a $lista$\;
	\Return $lista$\;
\end{algorithm}

\subsubsection{Esempio guidato}
Consideriamo il DAG:
\[
a\to b,\qquad a\to c,\qquad b\to d,\qquad c\to d.
\]
Un'esecuzione DFS puo' completare i nodi nell'ordine:
\[
d,\ b,\ c,\ a.
\]
Inserendo ogni nodo in testa alla lista al momento del completamento:
\[
\begin{array}{c|c}
\text{Nodo completato} & \text{Lista} \\
\toprule
d & \langle d\rangle \\
b & \langle b,d\rangle \\
c & \langle c,b,d\rangle \\
a & \langle a,c,b,d\rangle \\
\bottomrule
\end{array}
\]
La lista
\[
\langle a,c,b,d\rangle
\]
e' un ordinamento topologico: $a$ precede $b$ e $c$, e sia $b$ sia $c$ precedono $d$.

\subsubsection{Grafico o diagramma}
\begin{figure}[H]
	\centering
	\begin{tikzpicture}[
		v/.style={circle, draw=primary, thick, minimum size=7mm, font=\small},
		e/.style={-{Stealth[length=2mm]}, primary, thick}
		]
		\node[v] (a) at (0,1.4) {a};
		\node[v] (b) at (-1.1,0) {b};
		\node[v] (c) at (1.1,0) {c};
		\node[v] (d) at (0,-1.4) {d};
		\draw[e] (a)--(b) (a)--(c) (b)--(d) (c)--(d);
		\node[font=\scriptsize, below=3mm of d] {DAG: $\langle a,c,b,d\rangle$ è un ordinamento topologico};
	\end{tikzpicture}
	\caption{DAG per l'ordinamento topologico (architettura delle slide).}
\end{figure}

\begin{figure}[H]
	\centering
	\begin{tikzpicture}[
		v/.style={circle, draw=primary, thick, minimum size=7mm, font=\small},
		e/.style={-{Stealth[length=2mm]}, primary, thick}
		]
		\foreach \i/\lab in {0/a, 1/c, 2/b, 3/d}{
			\node[v] (o\i) at (\i*1.5,0) {\lab};
		}
		\draw[e] (o0)--(o1) (o0)--(o2) (o1)--(o3) (o2)--(o3);
		\node[font=\scriptsize, below=4mm of o1] {ordinamento $\langle a,c,b,d\rangle$: archi da sinistra a destra};
	\end{tikzpicture}
	\caption{Un ordinamento topologico valido disposto su una linea.}
\end{figure}

\subsubsection{Dimostrazione}
La correttezza si basa sul fatto indicato nelle slide:
\[
(u,v)\in E \Longrightarrow u.f>v.f
\]
in un DAG.

Consideriamo un arco $(u,v)$ durante DFS.

Se $v$ e' bianco quando l'arco viene esaminato, allora $v$ diventa discendente di $u$ nella foresta DFS. Un discendente viene completato prima dell'antenato, quindi
\[
v.f<u.f.
\]

Se $v$ e' nero, allora $v$ e' gia' stato completato prima del completamento di $u$, quindi ancora
\[
v.f<u.f.
\]

Il caso in cui $v$ sia grigio corrisponderebbe a un arco all'indietro verso un antenato. Questo produrrebbe un ciclo, impossibile in un DAG.

Quindi per ogni arco $(u,v)$ vale
\[
u.f>v.f.
\]
L'algoritmo inserisce ogni nodo in testa alla lista quando viene completato. Di conseguenza, i nodi con tempo di completamento maggiore compaiono prima. Pertanto, per ogni arco $(u,v)$, $u$ compare prima di $v$. Questa e' esattamente la definizione di ordinamento topologico.

\subsubsection{Analisi della complessita'}
L'algoritmo esegue DFS e inserisce ogni nodo una volta in lista. Quindi:
\[
\Theta(|V|+|E|)
\]
per la DFS, piu' $O(|V|)$ per gli inserimenti in testa.

La complessita' finale e':
\[
\Theta(|V|+|E|).
\]
Lo spazio aggiuntivo e' $O(|V|)$ per la lista e per gli attributi di DFS.

\subsubsection{Errori tipici d'esame}
\begin{hardnote}
Applicare l'ordinamento topologico a un grafo diretto con cicli. La definizione usata nelle slide richiede un DAG.
\end{hardnote}

\begin{hardnote}
Ordinare i nodi per tempo di scoperta invece che per tempo di completamento. L'algoritmo usa i tempi $f$, non i tempi $d$.
\end{hardnote}

\begin{hardnote}
Inserire il nodo in coda invece che in testa al completamento. L'inserimento in testa realizza l'ordine decrescente dei tempi di completamento.
\end{hardnote}

\subsubsection{Possibili domande d'esame}
\begin{itemize}
	\item Definire DAG e ordinamento topologico.
	\item Scrivere l'algoritmo \textsc{Topological-Sort}.
	\item Dimostrare che in un DAG, per ogni arco $(u,v)$, vale $u.f>v.f$.
	\item Simulare l'algoritmo su un piccolo DAG.
	\item Analizzare la complessita'.
\end{itemize}

\subsubsection{Collegamenti teorici}
L'ordinamento topologico e' un'applicazione diretta di DFS. Usa i tempi di completamento e il fatto che un DAG non contiene archi all'indietro.

% ===============================================================
\subsection{Componenti fortemente connesse}

\subsubsection{Definizione}
In un grafo diretto
\[
G=(V,E),
\]
una \textbf{componente fortemente connessa} e' un insieme massimale di nodi
\[
C\subseteq V
\]
tale che, per ogni coppia $u,v\in C$, vale:
\[
u\leadsto v
\qquad\text{e}\qquad
v\leadsto u.
\]

Il termine \textbf{massimale} significa che non e' possibile aggiungere a $C$ un altro nodo preservando la proprieta' di mutua raggiungibilita'.

\paragraph{Grafo trasposto.}
Dato un grafo diretto
\[
G=(V,E),
\]
il suo grafo trasposto e'
\[
G^T=(V,E^T),
\]
dove
\[
E^T=\{(v,u)\mid (u,v)\in E\}.
\]
Il grafo trasposto inverte il verso di tutti gli archi.

\subsubsection{Intuizione}
Una componente fortemente connessa e' una regione del grafo in cui tutti i nodi sono raggiungibili reciprocamente. Se si contraggono le componenti in singoli nodi, si ottiene una struttura senza cicli tra componenti.

L'algoritmo delle slide usa due visite DFS:
\begin{enumerate}
	\item una DFS su $G$ per calcolare i tempi di completamento;
	\item una DFS su $G^T$, scegliendo i nodi in ordine decrescente rispetto ai tempi di completamento calcolati in $G$.
\end{enumerate}

\subsubsection{Motivazione}
DFS da sola non basta a separare direttamente le componenti fortemente connesse in un grafo diretto arbitrario. I tempi di completamento della prima DFS indicano l'ordine corretto in cui esplorare il grafo trasposto.

\subsubsection{Proprieta' teoriche}
\paragraph{Stesse componenti in $G$ e $G^T$.}
I grafi $G$ e $G^T$ hanno le stesse componenti fortemente connesse. Infatti, se in $G$ esiste un cammino $u\leadsto v$, allora in $G^T$ esiste il cammino inverso $v\leadsto u$. La mutua raggiungibilita' e' quindi preservata invertendo tutti gli archi.

\paragraph{Tempi di componente.}
Per una componente $U$, definiamo:
\[
d(U)=\min\{u.d\mid u\in U\},
\]
cioe' il primo tempo di scoperta di un vertice della componente, e
\[
f(U)=\max\{u.f\mid u\in U\},
\]
cioe' l'ultimo tempo di completamento di un vertice della componente.

\begin{lemma}
Si considerino due componenti fortemente connesse distinte $C$ e $C'$ di $G=(V,E)$. Se
\[
(u,v)\in E,\qquad u\in C',\qquad v\in C,
\]
allora
\[
f(C')>f(C).
\]
\end{lemma}

\begin{corollario}
Si considerino due componenti fortemente connesse distinte $C$ e $C'$ di $G=(V,E)$. Se
\[
f(C')<f(C),
\]
allora $E^T$ non contiene nessun arco $(v,u)$ tale che
\[
v\in C,\qquad u\in C'.
\]
\end{corollario}

\subsubsection{Algoritmo}
\begin{algorithm}[H]
	\caption{\textsc{Strongly-Connected-Components}$(G)$}
	\KwIn{$G=(V,E)$}
	\KwOut{Componenti fortemente connesse di $G$}
	\textsc{DFS}$(G)$\tcp*{determina i tempi di completamento $u.f$}
	Determina $G^T$\;
	Esegui \textsc{DFS}$(G^T)$ esaminando i nodi in ordine decrescente rispetto al tempo di completamento calcolato in $G$\;
	$C\gets[\,]$\;
	$i\gets 0$\;
	\ForEach{albero DFS $T$ in $G^T$}{
		$C[i]\gets$ sottografo indotto dai nodi dell'albero $T$\;
		$i\gets i+1$\;
	}
	\Return $C$\;
\end{algorithm}

\subsubsection{Esempio guidato}
Consideriamo il grafo diretto:
\[
a\to b,\quad b\to a,\quad b\to c,\quad c\to d,\quad d\to c.
\]
Le componenti fortemente connesse sono:
\[
C_1=\{a,b\},\qquad C_2=\{c,d\}.
\]
Infatti $a$ e $b$ si raggiungono reciprocamente, e anche $c$ e $d$ si raggiungono reciprocamente. Inoltre esiste un arco da $C_1$ a $C_2$, ma non un cammino di ritorno da $C_2$ a $C_1$.

La prima DFS su $G$ calcola i tempi $f$. Poi si costruisce $G^T$, invertendo tutti gli archi:
\[
b\to a,\quad a\to b,\quad c\to b,\quad d\to c,\quad c\to d.
\]
La seconda DFS su $G^T$, usando l'ordine decrescente dei tempi $f$ della prima DFS, produce alberi DFS che coincidono con le componenti fortemente connesse.

\subsubsection{Grafico o diagramma}
\begin{figure}[H]
	\centering
	\begin{tikzpicture}[
		v/.style={circle, draw=primary, thick, minimum size=7mm, font=\small},
		e/.style={-{Stealth[length=2mm]}, primary, thick},
		comp/.style={draw=primary!50, dashed, rounded corners, inner sep=4pt}
		]
		\node[v] (a) at (0,0.8) {a};
		\node[v] (b) at (1.6,0.8) {b};
		\node[v] (c) at (3.4,0.8) {c};
		\node[v] (d) at (5.0,0.8) {d};
		\draw[e] (a) to[bend left=18] (b);
		\draw[e] (b) to[bend left=18] (a);
		\draw[e] (b)--(c);
		\draw[e] (c) to[bend left=18] (d);
		\draw[e] (d) to[bend left=18] (c);
		\node[comp, fit=(a)(b)] (cab) {};
		\node[comp, fit=(c)(d)] (ccd) {};
		\node[font=\scriptsize, below=2mm of cab] {$\{a,b\}$};
		\node[font=\scriptsize, below=2mm of ccd] {$\{c,d\}$};
	\end{tikzpicture}
	\caption{Grafo orientato $G$ e componenti fortemente connesse (esempio SCC delle slide).}
\end{figure}

\begin{figure}[H]
	\centering
	\begin{tikzpicture}[
		v/.style={circle, draw=primary, thick, minimum size=7mm, font=\small},
		e/.style={-{Stealth[length=2mm]}, primary, thick},
		comp/.style={draw=primary!50, dashed, rounded corners, inner sep=4pt},
		gt/.style={font=\small\bfseries, primary}
		]
		\node[gt] at (-0.3,1.55) {$G$};
		\node[v] (a) at (0,0.8) {a};
		\node[v] (b) at (1.5,0.8) {b};
		\node[v] (c) at (3.2,0.8) {c};
		\node[v] (d) at (4.7,0.8) {d};
		\draw[e] (a) to[bend left=18] (b);
		\draw[e] (b) to[bend left=18] (a);
		\draw[e] (b)--(c);
		\draw[e] (c) to[bend left=18] (d);
		\draw[e] (d) to[bend left=18] (c);
		\node[comp, fit=(a)(b)] {};
		\node[comp, fit=(c)(d)] {};
		\draw[-{Stealth[length=2.5mm]}, primary, thick] (2.35,0.35) -- (2.85,0.35);

		\node[gt] at (6.0,1.55) {$G^T$};
		\node[v] (at) at (6.0,0.8) {a};
		\node[v] (bt) at (7.5,0.8) {b};
		\node[v] (ct) at (9.2,0.8) {c};
		\node[v] (dt) at (10.7,0.8) {d};
		\draw[e] (at) to[bend left=18] (bt);
		\draw[e] (bt) to[bend left=18] (at);
		\draw[e] (ct)--(bt);
		\draw[e] (dt) to[bend left=18] (ct);
		\draw[e] (ct) to[bend left=18] (dt);
		\node[comp, fit=(at)(bt)] {};
		\node[comp, fit=(ct)(dt)] {};
	\end{tikzpicture}
	\caption{$G$, $G^T$ e componenti $\{a,b\}\to\{c,d\}$: la seconda DFS su $G^T$ estrae un albero per componente.}
\end{figure}

\subsubsection{Dimostrazione}
\paragraph{Dimostrazione del lemma.}
Siano $C$ e $C'$ due componenti fortemente connesse distinte e sia
\[
(u,v)\in E
\]
con
\[
u\in C',\qquad v\in C.
\]

Ci sono due casi.

\textbf{Caso 1: $d(C')<d(C)$.}
Sia $y\in C'$ tale che
\[
y.d=d(C').
\]
Al tempo $y.d$, i nodi di $C'$ e di $C$ rilevanti per il cammino sono ancora bianchi, e da $y$ si possono raggiungere:
\begin{itemize}
	\item tutti i nodi di $C'$, per forte connessione;
	\item il nodo $u$, quindi tramite l'arco $(u,v)$ il nodo $v$;
	\item tutti i nodi di $C$, per forte connessione di $C$.
\end{itemize}
Quindi i nodi di $C'$ e $C$ diventano discendenti di $y$ nella DFS. Ne segue:
\[
y.f=f(C')>f(C).
\]

\textbf{Caso 2: $d(C')>d(C)$.}
Sia $y\in C$ tale che
\[
y.d=d(C).
\]
All'inizio della visita da $y$, i nodi di $C'$ sono ancora bianchi. Non puo' esistere un cammino da $y$ a un nodo di $C'$, perche' altrimenti, usando l'arco da $C'$ a $C$ e la forte connessione interna delle due componenti, $C$ e $C'$ non sarebbero due componenti distinte.

Dunque la DFS partita da $y$ completa i nodi di $C$ senza scoprire i nodi di $C'$. Al tempo
\[
y.f=f(C),
\]
i nodi di $C'$ sono ancora bianchi. Essi saranno scoperti e completati solo dopo. Quindi:
\[
f(C')>f(C).
\]

In entrambi i casi vale
\[
f(C')>f(C).
\]

\paragraph{Dimostrazione del corollario.}
La contronominale del lemma dice: se
\[
f(C')<f(C),
\]
allora non esiste alcun arco
\[
(u,v)\in E
\]
con
\[
u\in C',\qquad v\in C.
\]
Nel grafo trasposto, un arco $(u,v)$ di $G$ diventa $(v,u)$ in $G^T$. Quindi $G^T$ non contiene archi da $C$ verso $C'$.

\paragraph{Correttezza dell'algoritmo SCC.}
La prima DFS su $G$ calcola i tempi di completamento. La seconda DFS viene eseguita su $G^T$ in ordine decrescente di tali tempi.

Sia $x$ il primo nodo scelto nella seconda DFS, cioe' quello con tempo $x.f$ massimo. Sia $C$ la componente fortemente connessa che contiene $x$. Allora $f(C)$ e' massimo tra le componenti.

Per il corollario, in $G^T$ non esistono archi da $C$ verso componenti con tempo di completamento minore ancora bianche. Quindi la visita DFS da $x$ in $G^T$ non puo' uscire da $C$ verso una nuova componente bianca. Poiche' $G$ e $G^T$ hanno le stesse componenti fortemente connesse, la DFS da $x$ raggiunge tutti e soli i nodi di $C$. Il primo albero DFS della seconda visita determina quindi una componente fortemente connessa.

Nel passo successivo, il ciclo principale di DFS su $G^T$ sceglie tra i nodi ancora bianchi un nodo $y$ con tempo di completamento massimo. Lo stesso ragionamento vale per la componente di $y$: eventuali archi verso componenti gia' trovate portano a nodi neri, mentre non esistono archi verso componenti bianche con tempo minore che possano far uscire la visita dalla componente corrente.

Ripetendo il ragionamento, ogni albero DFS della seconda visita identifica esattamente una componente fortemente connessa. Quindi l'algoritmo calcola correttamente tutte le componenti fortemente connesse di un grafo orientato.

\subsubsection{Analisi della complessita'}
L'algoritmo esegue:
\begin{itemize}
	\item una DFS su $G$:
	\[
	\Theta(|V|+|E|);
	\]
	\item la costruzione di $G^T$:
	\[
	\Theta(|V|+|E|)
	\]
	se $G$ e' rappresentato con liste di adiacenza;
	\item una DFS su $G^T$:
	\[
	\Theta(|V|+|E|).
	\]
\end{itemize}

La complessita' totale e':
\[
\Theta(|V|+|E|).
\]

Le slide osservano che ciascun $C[i]$ puo' essere determinato durante la seconda visita DFS; il ciclo che raccoglie gli alberi e' scritto nell'algoritmo per chiarezza.

\subsubsection{Errori tipici d'esame}
\begin{hardnote}
Eseguire la seconda DFS su $G$ invece che su $G^T$. L'inversione degli archi e' essenziale.
\end{hardnote}

\begin{hardnote}
Usare nella seconda DFS un ordine qualunque dei nodi. L'ordine deve essere decrescente rispetto ai tempi di completamento calcolati nella prima DFS su $G$.
\end{hardnote}

\begin{hardnote}
Confondere componente connessa e componente fortemente connessa. Nei grafi diretti serve raggiungibilita' in entrambe le direzioni per ogni coppia di nodi della componente.
\end{hardnote}

\subsubsection{Possibili domande d'esame}
\begin{itemize}
	\item Definire una componente fortemente connessa.
	\item Definire il grafo trasposto.
	\item Dimostrare che $G$ e $G^T$ hanno le stesse componenti fortemente connesse.
	\item Enunciare e dimostrare il lemma sui tempi $f(C')>f(C)$.
	\item Scrivere l'algoritmo \textsc{Strongly-Connected-Components}.
	\item Dimostrare la correttezza dell'algoritmo.
	\item Analizzare la complessita' $\Theta(|V|+|E|)$.
\end{itemize}

\subsubsection{Collegamenti teorici}
Le componenti fortemente connesse sono un'applicazione avanzata di DFS. L'algoritmo usa:
\begin{itemize}
	\item tempi di completamento;
	\item grafo trasposto;
	\item seconda DFS ordinata per tempi decrescenti.
\end{itemize}

% ===============================================================
\subsection{Procedura per stampare un cammino minimo}

\subsubsection{Definizione}
Tra gli esercizi ufficiali delle slide compare la richiesta di scrivere la procedura che stampa il cammino minimo dalla sorgente a un nodo dato, se esiste.

Dopo BFS, per ogni nodo raggiungibile $v$, il predecessore $v.\pi$ permette di risalire da $v$ verso la sorgente $s$ lungo l'albero BF.

\subsubsection{Intuizione}
Il cammino viene ricostruito all'indietro:
\[
v,\ v.\pi,\ v.\pi.\pi,\ \ldots,\ s.
\]
Per stamparlo nell'ordine corretto da $s$ a $v$, si usa una procedura ricorsiva: prima stampa il cammino fino al predecessore, poi stampa $v$.

\subsubsection{Motivazione}
BFS non restituisce solo le distanze. Gli attributi $\pi$ codificano anche un cammino minimo effettivo, come affermato dal teorema di correttezza.

\subsubsection{Proprieta' teoriche}
Se $v$ e' raggiungibile da $s$, allora la catena dei predecessori termina in $s$:
\[
v,\ v.\pi,\ v.\pi.\pi,\ldots,s.
\]
Se $v$ non e' raggiungibile, allora
\[
v.\pi=\textsc{Null}
\]
e $v\neq s$.

\subsubsection{Algoritmo}
\begin{algorithm}[H]
	\caption{\textsc{Print-Path}$(s,v)$}
	\KwIn{Sorgente $s$, nodo destinazione $v$, predecessori $\pi$ calcolati da BFS}
	\KwOut{Cammino minimo da $s$ a $v$, se esiste}
	\If{$v=s$}{
		stampa $s$\;
	}
	\Else{
		\If{$v.\pi=\textsc{Null}$}{
			stampa ``nessun cammino da $s$ a $v$''\;
		}
		\Else{
			\textsc{Print-Path}$(s,v.\pi)$\;
			stampa $v$\;
		}
	}
\end{algorithm}

\subsubsection{Esempio guidato}
Nell'esempio BFS con sorgente $2$:
\[
5.\pi=4,\qquad 4.\pi=2,\qquad 2.\pi=\textsc{Null}.
\]
La chiamata
\[
\textsc{Print-Path}(2,5)
\]
esegue:
\[
\textsc{Print-Path}(2,4),
\]
che esegue:
\[
\textsc{Print-Path}(2,2).
\]
Stampa quindi:
\[
2,\ 4,\ 5.
\]

\subsubsection{Grafico o diagramma}
\begin{figure}[H]
	\centering
	\begin{tikzpicture}[
		v/.style={circle, draw=primary, thick, minimum size=7mm, font=\small},
		e/.style={primary, thick},
		pi/.style={-{Stealth[length=2mm]}, primary, thick},
		lbl/.style={font=\scriptsize, primary}
		]
		\node[v] (n2) at (0,1.2) {2};
		\node[v] (n1) at (-1.2,0) {1};
		\node[v] (n4) at (0.4,0) {4};
		\node[v] (n5) at (1.8,-0.8) {5};
		\draw[e, gray!45] (n1)--(n2) (n1)--(n4) (n2)--(n4) (n4)--(n5);
		\draw[pi] (n5)--(n4) node[midway, right, lbl] {$\pi$};
		\draw[pi] (n4)--(n2);
		\node[lbl, above=1mm of n2] {$2.\pi=\textsc{Null}$};
		\node[font=\scriptsize, below=3mm of n5] {
			\textsc{Print-Path}$(2,5)$ stampa $2\to 4\to 5$
		};
	\end{tikzpicture}
	\caption{Catena dei predecessori da BFS per ricostruire il cammino minimo $2\to 4\to 5$.}
\end{figure}

\subsubsection{Dimostrazione}
La procedura e' corretta per induzione sulla lunghezza del cammino.

\textbf{Caso base.}
Se $v=s$, il cammino da $s$ a $s$ contiene solo $s$, e la procedura stampa $s$.

\textbf{Passo induttivo.}
Se $v\neq s$ e $v.\pi\neq\textsc{Null}$, per il teorema di correttezza di BFS un cammino minimo da $s$ a $v$ e' formato da un cammino minimo da $s$ a $v.\pi$ seguito dall'arco $(v.\pi,v)$. Per ipotesi induttiva, la chiamata ricorsiva stampa correttamente il cammino da $s$ a $v.\pi$; stampando poi $v$, si ottiene il cammino da $s$ a $v$.

Se $v.\pi=\textsc{Null}$ e $v\neq s$, non esiste cammino da $s$ a $v$ registrato da BFS, quindi la procedura segnala correttamente l'assenza del cammino.

\subsubsection{Analisi della complessita'}
La procedura segue al piu' un predecessore per ogni nodo del cammino. Nel caso peggiore:
\[
O(|V|)
\]
tempo e
\[
O(|V|)
\]
spazio sullo stack ricorsivo.

\subsubsection{Errori tipici d'esame}
\begin{hardnote}
Stampare prima $v$ e poi il predecessore produce il cammino al contrario. La procedura ricorsiva deve stampare prima il cammino verso $v.\pi$ e poi $v$.
\end{hardnote}

\begin{hardnote}
Non gestire il caso $v.\pi=\textsc{Null}$ quando $v\neq s$. In quel caso il cammino non esiste.
\end{hardnote}

\subsubsection{Possibili domande d'esame}
\begin{itemize}
	\item Scrivere \textsc{Print-Path}.
	\item Spiegare perche' usa i predecessori $\pi$.
	\item Dimostrare la correttezza per induzione.
	\item Analizzare tempo e spazio.
\end{itemize}

\subsubsection{Collegamenti teorici}
La procedura e' collegata direttamente al teorema di correttezza di BFS: le distanze $d$ danno la lunghezza minima, mentre i predecessori $\pi$ permettono di ricostruire un cammino minimo.
